\scanpage{9}\mspage{8}
$\inf$ in $\unsure{\mathcal{F}|\mathfrak{F}}$ is $\varnothing_{\mathcal{F}}$, the least element. of $\mathcal{F}$ or \uline{empty figure}. Thus
\[
F\text{ and }G\text{ disjoint}\quad\Longrightarrow\quad F\text{ and }G\text{ compatible}.
\]

We have, by the axioms:
\[
\begin{aligned}
&&\deleted{\text{every}}\quad\inserted{\operatorname{omb}(F)\text{ and }\operatorname{omb}(G)}\\
F\text{ and }G\text{ disjoint}&\quad\Longleftrightarrow\quad&\deleted{\operatorname{omb}(F)\cap\operatorname{omb}(G)\unsure{=}}\\
&&\text{are ``disjoint'' \unsure{in the|a} sense}\\
&&\text{that follows: they are \unsure{disjoint|disjoint}}\\
&&\text{set-theoretically i.e.}\\
&&\operatorname{omb}(F)\cap\operatorname{om}(G)=\varnothing,\ \text{and}\\
&&\text{every element of }\operatorname{omb}(F)\text{ is}\\
&&\text{compatible with every element}\\
&&\text{of }\operatorname{om}(G)
\end{aligned}\tag*{(1.13)}
\]
-- i.e. two figures are disjoint if and only if they have no \unsure{common|common} locus, and every locus of one is \deleted{disjoint} \inserted{\deleted{com} disjoint} (in the sense of th\unsure{\deleted{e}|e}\inserted{'}\deleted{theory}workshop) from every locus of the other;

This shows that the relations of \deleted{\unsure{dis|disj}}compatibility and of disjointness are important also in the \deleted{\unsure{in some|in of}} \inserted{sub-}sets $\mathcal{M},\mathcal{L}$ of $\mathcal{F}$. In the case of set-theoretic worksh\unsure{ops|\deleted{op}ops|\deleted{ops}ops}, these relations are ``trivial'' on $\mathcal{L}$: two loci are always \deleted{disjoint} \inserted{compatible}, and disjoint if and only if they are distinct. On the other hand, even in the set-theoretic case and $\overline m$ tame, the relation
